\section{Introduction}

A self-avoiding path crossing a wide strip can make long excursions parallel
to its walls.  We prove a mean-length exponent of $4/3$ for critical honeycomb walks,
both with these excursions unrestricted and with their lateral displacement
bounded by a sufficiently large fixed multiple of the strip width.
The starting boundary point is fixed, the terminal point is summed over the
opposite wall, and all finite lengths receive their critical weights.  We
also obtain the corresponding mean for large half-plane arches and a bulk
visit law for chords in a triangle.

\subsection{The lattice and the positive path measures}
The honeycomb lattice has unit edges in directions
\[
 e_2=(0,-1),\qquad e_0=(\sqrt3/2,1/2),\qquad
 e_1=(-\sqrt3/2,1/2).
\]
Put $U=e_0-e_1$ and $V=e_0-e_2$.  The two vertex classes are
$jU+kV$ and $jU+kV+e_0$, for $j,k\in\mathbb Z$; the three edges from a
vertex of the first class have directions $e_0,e_1,e_2$.
For $i\in\{0,1,2\}$ and $t\in\frac12+\frac32\mathbb Z$, the line
$e_i\cdot y=t$ is an admissible wall.  It avoids vertices and cuts edges
orthogonally at their midpoints.  These midpoints are boundary \emph{ports}.  A port-to-port path visits distinct lattice vertices, all
strictly inside its prescribed domain.  Its two terminal half-edges have
no additional weight.

Throughout, a path $\gamma$ with $\ell(\gamma)$ visited vertices has weight
\[
 w(\gamma)=x^{\ell(\gamma)},\qquad
 x=(2\cos(\pi/8))^{-1}=(2+\sqrt2)^{-1/2}.
\]
The connective constant of the honeycomb lattice is $x^{-1}$ by
Duminil-Copin and Smirnov~\cite[Theorem~1]{DuminilCopinSmirnov2012}.
For a nonempty path between ordinary vertices, vertex count is edge count
plus one.  All normalizations below use vertex count.
We write $f(R)\asymp g(R)$ when their ratio stays between positive constants,
and $f(R)=R^{a+o(1)}$ when $\log f(R)/\log R\to a$.

Fix two parallel admissible walls at distance $R\in\frac32\mathbb Z_{>0}$,
call them the lower and upper walls, and fix a port $a$ on the lower wall.
Let $\tau$ be either unit tangent to the walls and write
$\xi(y)=\tau\cdot(y-a)$ for lateral displacement from the source.
A \emph{bridge} joins $a$ to any port of the upper wall and otherwise stays
strictly between the walls.  Write
\[
 B(R)=\sum_{\gamma\text{ bridge}}w(\gamma),\qquad
 M(R)=\sum_{\gamma\text{ bridge}}w(\gamma)\ell(\gamma).
\]
The normalized critical bridge law assigns probability $w(\gamma)/B(R)$.
For $C>0$, let $B_C(R)$ and $M_C(R)$ be the same sums restricted to paths
satisfying $\sup_{y\in\gamma}|\xi(y)|\le CR$, where the supremum includes
the terminal ports.  The sign chosen for $\tau$ does not affect this restriction.

\begin{theorem}[Bridge mass retained in a fixed corridor]
\label{thm:confined-bridge}
There is a constant $C_0<\infty$ such that, for every fixed $C\ge C_0$,
\[
 B_C(R)\asymp B(R)\asymp R^{-1/4},\qquad
 M_C(R)=R^{13/12+o(1)},\qquad M(R)=R^{13/12+o(1)}.
\]
Consequently the mean number of visited vertices is $R^{4/3+o(1)}$
both under the unrestricted critical bridge law and under its restriction
to this fixed corridor.
\end{theorem}

The constant $C_0$ is independent of the exponent slack implicit in $o(1)$.
The theorem concerns a mean under a law that sums all lengths.  Its proof
will expose why two further distinctions matter: a bulk visit estimate
initially counts both arches and bridges, and an endpoint-separation
estimate is needed before two disjoint exterior arms can turn those visits
into confined bridges.

\subsection{Boundary and finite-domain consequences}
A half-plane \emph{arch} starts at one fixed port and ends at another port
of the same wall.  Let $h(m)$ be its total critical mass when the terminal
port lies $m>0$ port spacings to the right.  Let $\mu$ be the unnormalized
critical measure with that initial port and the terminal port summed over
the wall.  Height is distance from the wall; diameter is Euclidean diameter.

\begin{theorem}[Sharp boundary mass and annealed chord laws]
\label{thm:boundary}\label{c2:thm-annealed}
The following conclusions hold as the indicated scale tends to infinity.
\begin{enumerate}
\item The arch kernel is nonincreasing and
\[
 h(m)\asymp m^{-5/4},\qquad
 \mu(\operatorname{height}>R)\asymp
 \mu(\operatorname{diameter}>R)\asymp R^{-1/4}.
\]
\item For every sufficiently large fixed $C$, normalize $\mu$ either on
$R\le\operatorname{diameter}\le CR$, or on
$\operatorname{height}\ge R$ and $\operatorname{diameter}\le CR$.
Under either law, $\mathbb E\ell=R^{4/3+o(1)}$.
\item Let
\[
 D_R=\{y:e_i\cdot y<R,\ i=0,1,2\},\qquad
 I_R=\{v\text{ a lattice vertex}:e_i\cdot v<(1-\varepsilon)R,
 \ i=0,1,2\},
\]
where $R\in\frac12+\frac32\mathbb Z$ and $\varepsilon>0$ is sufficiently
small and fixed.  Sum critical chord weights over the ordered pair of
boundary ports and normalize on chords visiting $I_R$.
Uniformly for $v\in I_R$, their visit probability is $R^{-2/3+o(1)}$.
The expected total number of visited vertices, and the expected number in
$I_R$, are both $R^{4/3+o(1)}$.
\end{enumerate}
\end{theorem}

The endpoints in the last two conclusions are summed.  Prescribing both
endpoints is a different conditioning problem.  The finite bridge reward
in Theorem~\ref{thm:confined-bridge} also supplies an input to renewal
arguments for infinite walks; those arguments concern different probability
laws and are not needed here.

\subsection{Historical context and related representations}
Nienhuis's solid-on-solid and Coulomb-gas analysis of the dilute $O(n)$
model gives the two-dimensional polymer exponent predictions at $n=0$
\cite{Nienhuis1982}. Lawler, Schramm and Werner formulated the conjectural
$\mathrm{SLE}_{8/3}$ scaling description of planar self-avoiding walk and
recovered the predicted displacement exponent $\nu=3/4$
\cite{LawlerSchrammWerner2004}. That prediction concerns
displacement under a fixed-length law. Our theorem determines a
mean length under the critical measure that sums lengths and terminal ports.
On the square lattice, Dyhr, Gilbert, Kennedy, Lawler and Passon studied
this critical strip ensemble with one fixed endpoint and the other summed,
relating it to a conditioned infinite half-plane walk \cite{DyhrEtAl2011}.

Duminil-Copin and Smirnov proved the exact honeycomb connective constant
using a para\-fermionic observable \cite{DuminilCopinSmirnov2012}.
Their strip argument gives $c/R\le B(R)\le1$ and records the sharper
power $R^{-1/4}$ as a prediction.
Beaton, Bousquet-M\'elou, de Gier, Duminil-Copin and Guttmann later proved
$B(R)\to0$ in their study of surface adsorption \cite{BeatonEtAl2014}.
Glazman and Manolescu gave a shorter proof and a logarithmic subsequence
bound, together with boundary two-point-function invariance for a
specified class of columnwise rhombic half-plane tilings
\cite{GlazmanManolescu2017}. Krachun and Panagiotis obtained a polynomial
upper bound on the critical bridge mass and quantitative sub-ballisticity
for uniform honeycomb walks \cite{KrachunPanagiotis2026}.

Our boundary exploration is a nonbacktracking form of the same
parafermionic cancellation, with winding phases and positive masses
kept distinct. Its relation to integrable rhombic loop weights was
developed by Ikhlef and Cardy \cite{IkhlefCardy2009}; Glazman gives
criticality and corrected general weight formulas in that setting
\cite{Glazman2015}. The local dilute diagram algebra and its
Yang--Baxter parametrization belong to the braid--monoid construction
of Grimm and Pearce \cite{GrimmPearce1993}. Zhou and Batchelor's
periodic transfer analysis supplies the spectral context, including a
seam controlling winding-loop weights \cite{ZhouBatchelor1997}.
We fix the boundary gauge, caps and scalar factors explicitly here.

Polynomial exchange equations and special-rapidity reductions have
counterparts in the dense $O(1)$ sum rule of Di Francesco and Zinn-Justin
\cite{DiFrancescoZinnJustin2005} and the dilute open-boundary $O(1)$
ground states and first fusion sum rule of Garbali and Nienhuis
\cite{GarbaliNienhuis2017Ground,GarbaliNienhuis2017Sum}.
Garbali and Nienhuis also observed that polynomial ground states and
the associated exchange and recurrence construction persist at $n=0$
\cite[Section~4]{GarbaliNienhuis2017Ground}.
Our $n=0$ disk vector, two-site compatibility and scalar normalization
are proved below. The reflected current calculation uses the classical
level-one lattice realization of Frenkel--Kac and Segal
\cite{FrenkelKac1980,Segal1981,Frenkel1985}; we give the mode and
vacuum-sign calculations required by the physical normalization.

The related polynomial-vacuum construction
\cite[Theorem~1.1]{PolynomialVacuumCompanion} gives the bridge mean exponent and
retains it in a corridor whose width is of order $R\log^2R$. The present argument gives
sharp boundary mass comparisons and confinement to a fixed multiple
of $R$. The two constructions have the same even scalar polynomial and,
after a change of lattice scale, the same homogeneous separator
partition function (Remark~\ref{rem:pd:scalar-comparison}).
The disk interpolation, normalization estimates and spatial arguments
used here are proved directly.

\subsection{Two exact calculations and the spatial conversion}
Our proof uses the dilute-loop representation at loop weight zero. One
calculation determines a boundary deficit. On a cylinder with $N$ sites
around its circumference, let $h_N(m)$ be the mass of a single upper-half-
cylinder arch from a fixed site to the site $m$ spacings to its right,
with that oriented boundary interval specifying its lift. We prove
\[
 1-2\cos(3\pi/8)\sum_{m=1}^{N-1}h_N(m)\asymp N^{-1/4}.
\]
We first deform the transfer weights by distinct rapidities. Their finite
permutation orbit lets us solve the endpoint equations as a matrix-inverse
pairing. At coincident rapidities this pairing is comparable to the squared norm
of polynomial evaluation against a positive measure. Its asymptotics give the deficit;
boundary exploration identities then convert it into height, strip and
pointwise endpoint bounds.

The second calculation concerns bulk polygons. Fix two opposite face
centers on a cylinder of circumference $2n|U|$, on the same horizontal cut.
Give each collection of pairwise vertex-disjoint simple polygons
separating the marks the product of weights $2x^{|\Gamma|}$; polygons
are unrooted and unoriented, and the empty collection has weight one.
We prove that this partition function is $n^{1/6+o(1)}$.
A finite disk transfer vector, normalized by polynomial interpolation,
gives the partition function as a pairing of two disk vectors.
To estimate it, replace the factor two per polygon by $2\cos\theta$.
The pairing has a scalar-gas representation. A reflection of its currents
gives an upper bound for complex $\theta$, while an independent determinant
calculation fixes the scalar normalization. The resulting upper exponent
is harmonic in $\theta$. At zero polygon fugacity the partition function
equals one. A maximum-principle argument, using this exact value and the
monotonicity of the positive polygon sum, forces the upper exponent to be
attained also at fugacity two.

To localize this cylinder calculation, we bound translation classes of
large polygons and winding polygons on a cylinder. The small polygons
around the two marks then contribute two copies of the same planar
nesting sum; the remaining polygons cost only a bounded factor. This gives
\[
 Z_R(f)=R^{1/12+o(1)},
\]
where $Z_R(f)$ sums nested polygons surrounding a planar face center $f$
and contained in its radius-$R$ ball. An exact finite exploration identity
converts this sum into the weight of boundary paths visiting a bulk vertex.
It gives the triangle law and, after sewing one exterior arm, the arch mean.

For bridges, both ends need exterior arms. A marked three-leg exploration
identity first shows that enough bulk-visit mass has well-separated chord
endpoints. Straight tubes with fixed endpoints, widening first tubes and
routes with at most two bends supply two disjoint arms. Summing over
possible sewing cuts produces many representations of one path; a weighted
second-moment estimate controls that multiplicity and recovers the mass of
distinct bridges. The resulting bridges stay in one fixed corridor.
Summing over order $R^2$ bulk vertices and accounting for order $R$
starting ports gives the raw first-moment exponent $1+1/12=13/12$.
The same nesting identity gives the unrestricted upper bound. Division
by the bridge mass then yields the mean exponent $4/3$.

The boundary calculation occupies Sections~\ref{c2:boundary}--\ref{sec:deficit}.
Sections~\ref{c2:sec-disk-gluing}--\ref{c2:sec-normalization}
prove the cylinder partition estimate.
Sections~\ref{c2:two-bridge}--\ref{c2:nesting} pass to planar nesting and
bulk visits. Section~\ref{c2:arches} proves the arch law;
Sections~\ref{c2:separated} and~\ref{c2:strips} establish endpoint separation
and complete the confined bridge proof.
